\section{Reference numerical method}\label{sec:methods}

Reference solutions are computed with solids4foam \citep{Cardiff2018}
\TODO{version / commit} on OpenFOAM \TODO{v2312/v2412}. Every setting in
this section is given in full, and is also contained in the released case
files, so that the solution can be reproduced exactly.
\Cref{tab:settings} collects all settings of the baseline configuration;
\cref{sec:ver:robust} reports how the solution and its robustness change
when each one is varied.

% ---------------------------------------------------------------------------
\subsection{Solid solver}\label{sec:methods:solid}
\paragraph{Formulation} Cell-centred finite volume, total Lagrangian,
total-displacement formulation
(\texttt{nonLinearGeometryTotalLagrangianTotalDisplacement}), material law
\texttt{viscoActiveHolzapfelOgden}.
\paragraph{Incompressibility} Mixed displacement--pressure formulation
(\texttt{solvePressure true}): the volumetric term
$\partial U/\partial J$ of the penalty energy is replaced by a cell-centred
pressure $p_{\mathrm s}$ that satisfies the same penalty relation, so the
material response is unchanged. The pure-displacement formulation is
compared in \cref{sec:ver:robust}.
\paragraph{Stabilisation} Collocated FV discretisations of the momentum
equation and of the pressure equation require stabilisation to suppress
checkerboard modes. Momentum: \texttt{diffStencilLaplacian}, i.e.\ the
difference between a compact and a wide-stencil Laplacian of $\uu$ scaled by
an implicit stiffness and a factor $c_{\uu}$. Pressure: Rhie--Chow with
factor $c_p$ (and $c_p^{\mathrm J}$ in the Jacobian).
\TODO{Write the stabilisation terms explicitly; state which stiffness scales
them (shear modulus vs.\ $\kappa$) and why this matters for a material with
$a=\SI{59}{Pa}$ and $\kappa=\SI{1e6}{Pa}$.}
\paragraph{Non-linear and linear solvers} Newton's method via PETSc SNES with
a matrix-free Jacobian and backtracking line search; \TODO{KSP/PC type,
absolute/relative tolerances, maximum iterations}.
\paragraph{Time integration} Second-order backward differentiation (BDF2)
for $\ddot\uu$ and $\dot\uu$; the viscous term uses $\dot\EE$ from the same
scheme. \TODO{Initialisation of BDF2 at the first step.}

% ---------------------------------------------------------------------------
\subsection{Fluid solver}\label{sec:methods:fluid}
PIMPLE pressure--velocity coupling on a moving mesh, \TODO{$n$ outer
correctors, $m$ pressure correctors, non-orthogonal correctors}, with the
Rhie--Chow-type face-flux correction for the time derivative
(\texttt{ddtCorr}). Mesh motion: \TODO{velocity Laplacian motion solver,
diffusivity (e.g.\ inverse distance from $\Gfsi$ or quadratic)}; the
discrete geometric conservation law is satisfied by computing face fluxes
from swept volumes. Schemes are listed in \cref{tab:settings}.

% ---------------------------------------------------------------------------
\subsection{Fluid--solid coupling}\label{sec:methods:coupling}
\paragraph{Algorithm} Partitioned and strongly coupled: within each time
step the fluid and solid solvers are iterated until the interface residual
\begin{equation}
r^{k} = \frac{\|\uu_{\Gfsi}^{k+1}-\uu_{\Gfsi}^{k}\|_{2}}
             {\max_{\Gfsi}\|\uu_{\Gfsi}^{k+1}\|}
\end{equation}
\TODO{check the exact residual definition used by solids4foam} falls below
$\epsilon_{\mathrm{FSI}}$ or $N_{\mathrm{FSI}}$ iterations are reached.
\paragraph{Interface conditions} The fluid receives the interface velocity
(Dirichlet) and the solid receives the fluid traction (Neumann). To
counter the added-mass effect, which is strong here because
$\rho_{\mathrm f}/\rho_{\mathrm s}\approx1.05$ \citep{Causin2005}, the
interface pressure is imposed through a Robin-type condition
(\texttt{elasticWallPressure}, \citealp{Tukovic2018}) with reference
pressure $p_{\mathrm{ref}}$ and reference velocity $U_{\mathrm{ref}}$.
\TODO{Give the Robin condition explicitly and the meaning of
\texttt{robinPressureReference}, \texttt{robinFluxReferenceVelocity}, and of
the \texttt{hsModel secant} option.}
\paragraph{Acceleration} Fixed under-relaxation $\omega$, Aitken dynamic
relaxation \citep{Kuttler2008} or IQN-ILS \citep{Degroote2009};
\TODO{baseline choice}.
\paragraph{Predictor} \TODO{Interface displacement extrapolated from previous
time steps (order?)}.

% ---------------------------------------------------------------------------
\subsection{Meshes}
Solid: the all-hexahedral O-grid meshes of the solids4foam AR25 tutorial,
levels M1--M4 with refinement ratio $r=2$ in each direction. Fluid:
\TODO{conforming hexahedral cavity mesh, nodes matching the solid at $\Gfsi$
or non-matching with interpolation; state the interpolation}.
\Cref{tab:meshes} lists the cell counts and quality metrics.

\begin{table}[t]
\centering
\caption{Mesh levels. Solid divisions: square side $\times$ ring $\times$ wall.}
\label{tab:meshes}
\small
\begin{tabular}{lrrrrr}
\toprule
Level & Solid divisions & Solid cells & Fluid cells & $h_{\Gfsi}$ [mm] & max non-orth. [\si{\degree}] \\
\midrule
M1 & $8\times18\times3$   & \num{1920}   & \TODO{} & \TODO{} & \TODO{} \\
M2 & $16\times36\times5$  & \num{12800}  & \TODO{} & \TODO{} & \TODO{} \\
M3 & $32\times72\times10$ & \num{102400} & \TODO{} & \TODO{} & \TODO{} \\
M4 & $64\times144\times20$ & \num{819200} & \TODO{} & \TODO{} & \TODO{} \\
\bottomrule
\end{tabular}
\end{table}
\TODO{M1$\to$M2 wall divisions go 3$\to$5 (ratio not exactly 2); use the
generalised GCI with an effective $r$ \citep{Celik2008}.}

% ---------------------------------------------------------------------------
\begin{table}[p]
\centering
\caption{Complete numerical settings of the baseline configuration
(OpenFOAM/solids4foam keywords in typewriter font). Starred entries are
varied in \cref{sec:ver:robust}.}
\label{tab:settings}
\scriptsize
\begin{tabularx}{\linewidth}{lX}
\toprule
\multicolumn{2}{l}{\textbf{Solid}} \\
Solid model & \texttt{nonLinearGeometryTotalLagrangianTotalDisplacement} \\
Incompressibility$^*$ & mixed $\uu$--$p$ (\texttt{solvePressure true}) \\
Momentum stabilisation$^*$ & \texttt{diffStencilLaplacian}, $c_{\uu}=100$ \\
Pressure stabilisation$^*$ & Rhie--Chow, $c_p = 1$, $c_p^{\mathrm J}=1$ \\
Gradient & \texttt{leastSquaresS4f} \\
Laplacian / snGrad & Gauss linear corrected / corrected \\
Time scheme ($\ddot\uu$, $\dot\uu$)$^*$ & \texttt{backward} (BDF2) \\
Non-linear solver & PETSc SNES Newton, matrix-free Jacobian, backtracking line search \\
Linear solver & \TODO{KSP / PC} \\
Tolerances & \TODO{SNES abs/rel}, \TODO{KSP rel} \\
\midrule
\multicolumn{2}{l}{\textbf{Fluid}} \\
Algorithm & PIMPLE, \texttt{ddtCorr true}, \TODO{nOuter, nCorr, nNonOrthCorr} \\
Time scheme$^*$ & \TODO{Euler (current) $\to$ backward} \\
Gradient & Gauss linear \\
Convection$^*$ & Gauss linearUpwind grad(U) \\
Laplacian / snGrad & Gauss linear corrected / corrected \\
Linear solvers & \TODO{p: GAMG, tol; U: smoothSolver/PBiCGStab, tol} \\
Mesh motion$^*$ & \TODO{solver, diffusivity} \\
\midrule
\multicolumn{2}{l}{\textbf{Coupling}} \\
Scheme & partitioned, strongly coupled \\
Interface pressure & \texttt{elasticWallPressure} (Robin), \texttt{hsModel secant} \\
Robin parameters$^*$ & $p_{\mathrm{ref}}=0.1$, $U_{\mathrm{ref}}=\SI{0.01}{m/s}$ \TODO{units/meaning} \\
Acceleration$^*$ & \texttt{fixedRelaxation}, $\omega = 1.0$ \\
Tolerance$^*$ & $\epsilon_{\mathrm{FSI}} = \num{5e-5}$ \\
Max.\ iterations & $N_{\mathrm{FSI}} = 30$ \\
\midrule
\multicolumn{2}{l}{\textbf{Time}} \\
Time step$^*$ & $\dt = \TODO{}$ (current tutorial: \SI{5e-5}{s}) \\
End time & \SI{1}{s} \\
\bottomrule
\end{tabularx}
\end{table}
